% jonatanbb.xyz — the CV, English. A terse print copy of the site: same words where it
% counts, same palette (ink, gold, light), none of the lighting. Self-contained: plain
% LaTeX + TeX Gyre Pagella (the Palatino the site's --serif stack asks for first).
% Build with ./build.sh — it installs the pdf as ../cv-en.pdf, the asset the site serves.
\documentclass[10pt,a4paper]{article}

\usepackage[T1]{fontenc}
\usepackage[utf8]{inputenc}
\usepackage{tgpagella}
\usepackage[a4paper,margin=1.45cm,bottom=1.1cm]{geometry}
\usepackage{graphicx}
\usepackage{xcolor}
\usepackage{parskip}
\usepackage[hidelinks]{hyperref}

% the site's palette (styles.css): warm near-black ink, the site gold, and a
% deepened gold for small print — pure gold hairlines wash out on paper
\definecolor{ink}{HTML}{1C1A16}
\definecolor{gold}{HTML}{FAD02A}
\definecolor{goldink}{HTML}{A8890F}
\definecolor{soft}{HTML}{55514A}
\color{ink}
\hypersetup{colorlinks=true, urlcolor=goldink, linkcolor=goldink}

\pagestyle{empty}
\setlength{\parskip}{0pt}

% section header, echoing the site's kicker + bar: the name in gold serif
% lowercase, then the thick black rule running to the margin
\newcommand{\cvsection}[1]{%
  \par\vspace{8pt}%
  \noindent{\large\color{goldink}#1}\hspace{12pt}%
  \raisebox{3pt}{\color{ink}\leaders\hrule height 2.2pt\hfill\kern0pt}%
  \par\vspace{5pt}}

% timeline entry, like the site's: dates in a left column, the what on the right
\newcommand{\entry}[4]{%
  \noindent
  \begin{minipage}[t]{2.6cm}\raggedright\small\textbf{#1}\end{minipage}%
  \hspace{.45cm}%
  \begin{minipage}[t]{\dimexpr\textwidth-3.05cm\relax}
    \textbf{#2}\\[1pt]
    {\small\color{soft}#3}#4
  \end{minipage}\par\vspace{6.5pt}}
\newcommand{\entrybody}[1]{\\[2.5pt]{\small #1}}

\begin{document}

% ---- header: name + eyebrow + contact on the left, the framed portrait right ----
\noindent
\begin{minipage}[b]{\dimexpr\textwidth-3.6cm\relax}
  {\Huge Jonatan Barreiro Bastos}\\[7pt]
  {\small\scshape mathematician \textperiodcentered\ predoctoral researcher}\\[9pt]
  {\small O Porriño, Pontevedra, Spain \,\textperiodcentered\, +34 678 25 50 96 \,\textperiodcentered\, \href{mailto:jonatanbb@tutanota.com}{jonatanbb@tutanota.com}}\\[2pt]
  {\small\href{https://jonatanbb.xyz}{jonatanbb.xyz} \,\textperiodcentered\, Telegram \href{https://t.me/jonatanbb}{@jonatanbb} \,\textperiodcentered\, ORCID \href{https://orcid.org/0000-0003-0922-8055}{0000-0003-0922-8055}}
\end{minipage}\hfill
\begin{minipage}[b]{3.1cm}
  \raggedleft{\setlength{\fboxsep}{0pt}\setlength{\fboxrule}{1.1pt}\fcolorbox{ink}{white}{\includegraphics[width=2.9cm]{portrait.jpg}}}
\end{minipage}
\par\vspace{8pt}
{\color{gold}\hrule height 2.2pt}
\vspace{9pt}

% ---- profile: the site's lede, terse ----
\noindent Mathematician by training, of the mostly theoretical kind — complex, functional
and harmonic analysis — with a lasting knack for modelling problems, mathematical or not.
Four years of applied research on LiDAR point clouds added the practical layer: C++,
Python, and the everyday craft of building software. Now finishing my PhD and looking to
move from academia into industry.

\cvsection{experience}

\entry{Apr 2021 --\\Jul 2025}
  {Senior Research Technician}
  {CiTIUS — Centro de Investigación en Tecnoloxías Intelixentes, USC}
  {\entrybody{Predoctoral researcher in the LiDAR group under Dr.~Francisco F. Rivera.
   Modelled and implemented detection and integration techniques for LiDAR and image
   point clouds, primarily in C++. Since August 2025, continuing the underlying doctoral
   research independently.}}

\entry{2018 -- 2019}
  {Graduate Teaching Assistant}
  {University of Kansas}
  {\entrybody{Winter semester — instructor of record for a 33-student section of Math 115
   Calculus I, coordinated by Dr.~W. Huang and A. Kolasinski. Lectured, designed and
   graded the tests, ran the Help Room.\\[2pt]
   Spring semester — teaching assistant for Math 147 Calculus III (Honors), with
   Prof.~E. Gavosto: the highest calculus course graduate students were entrusted with.
   Taught classes, graded, ran the Help Room, and handled the 3D printing of student
   projects.}}

\entry{2017 -- 2018}
  {Graduate Research Assistant}
  {University of Kansas}
  {\entrybody{Research assistant under Prof.~Rodolfo Torres, beginning research in
   harmonic analysis.}}

\cvsection{education}

\entry{2017 -- 2020}
  {Master of Arts in Mathematics}
  {University of Kansas, USA \,\textperiodcentered\, GPA 4.0 / 4.0}
  {\entrybody{Recipient of the Dean's Doctoral Fellowship.}}

\entry{2015 -- 2016}
  {Máster en Matemáticas y Aplicaciones}
  {Universidad Autónoma de Madrid \,\textperiodcentered\, 8.0 / 10.0}
  {\entrybody{Research initiation track. Thesis: \emph{El Teorema de la Corona} (The
   Corona Theorem), advised by J. L. Fernández.}}

\entry{2011 -- 2015}
  {Grado en Matemáticas}
  {Universidad Autónoma de Madrid \,\textperiodcentered\, 7.6 / 10.0}
  {\entrybody{Thesis: \emph{La Transformada Rápida de Fourier y Aplicaciones} (The Fast
   Fourier Transform and Applications), advised by E. Hernández.}}

\cvsection{research}

\entry{In writing}
  {A Decoupled Random-Sampling Hough Approach to Line Detection in 2D Point Clouds}
  {Research complete \,\textperiodcentered\, extends to circle detection \,\textperiodcentered\, aimed at the IPOL journal}
  {}

\entry{Feb 2026}
  {LitS: A novel Neighbourhood Descriptor for Point Clouds}
  {arXiv:\href{https://arxiv.org/abs/2602.04838}{2602.04838} \,\textperiodcentered\, with F. F. Rivera, O. G. Lorenzo, D. L. Vilariño, J. C. Cabaleiro, A. M. Esmorís and T. F. Pena}
  {}

\cvsection{skills}

\entry{Programming}
  {\normalfont\small Proficient: \LaTeX, Sage, C++, Python \,\textperiodcentered\, Somewhat familiar: C, Java, MATLAB, R}
  {}{}

\entry{Languages}
  {\normalfont\small Spanish and Galician (native) \,\textperiodcentered\, English (fluent)}
  {}{}

\vfill
\begin{center}
  \small\color{soft}Updated July 13, 2026 \,\textperiodcentered\, the long version lives at \href{https://jonatanbb.xyz}{jonatanbb.xyz}
\end{center}

\end{document}
